% =====================================================================
%  chapters/minggu-03.tex
%  Minggu 3 — HTML Dasar: Struktur & Semantic Tag
% =====================================================================

\chapter{HTML Dasar: Struktur \& Semantic Tag}
\label{chap:minggu-03}

% ---------------------------------------------------------------------
% 1. CAPAIAN PEMBELAJARAN
% ---------------------------------------------------------------------
\begin{capaian}
Setelah menyelesaikan bab ini, mahasiswa diharapkan mampu:
\begin{itemize}
  \item menjelaskan struktur dasar dokumen HTML (\emph{doctype}, \texttt{html}, \texttt{head}, \texttt{body});
  \item menggunakan elemen HTML dasar: \emph{heading}, paragraf, tautan, gambar, dan daftar;
  \item membedakan dan menggunakan \emph{semantic tag} (\texttt{<header>}, \texttt{<nav>}, \texttt{<main>}, \texttt{<section>}, \texttt{<footer>}, dsb.);
  \item menulis HTML yang valid menggunakan W3C Validator;
  \item menghubungkan file HTML ke file CSS eksternal secara dasar.
\end{itemize}
\end{capaian}

% ---------------------------------------------------------------------
% 2. TUJUAN PRAKTIK
% ---------------------------------------------------------------------
\begin{tujuanpraktik}
Pada sesi praktik ini, mahasiswa akan:
\begin{itemize}
  \item membuat dokumen HTML dengan struktur yang benar;
  \item berlatih menggunakan elemen dasar HTML;
  \item membangun halaman profil dengan \emph{semantic tag};
  \item menghubungkan file HTML ke file CSS;
  \item memvalidasi kode HTML melalui W3C Validator.
\end{itemize}
\end{tujuanpraktik}

% ---------------------------------------------------------------------
% 3. STRUKTUR DOKUMEN HTML
% ---------------------------------------------------------------------
\section{Struktur Dokumen HTML}
\label{sec:m03-struktur}

\subsection{Proyek Baru}

\begin{langkah}
  \item Di VS Code, buat folder baru: \textbf{File\,$\rightarrow$\,Open Folder\,$\rightarrow$\,New Folder}, beri nama \texttt{latihan-web-3}.
  \item Buat file baru: \texttt{index.html}.
\end{langkah}

\subsection{Struktur Dasar}

Ketik kode berikut di \texttt{index.html}:

\begin{lstlisting}[language=HTML, caption={Struktur dasar dokumen HTML}, label={lst:m03-struktur}]
<!DOCTYPE html>
<html lang="id">
<head>
    <meta charset="UTF-8">
    <meta name="viewport" content="width=device-width, initial-scale=1.0">
    <title>Latihan HTML --- Minggu 3</title>
</head>
<body>
    <h1>Halo dari HTML!</h1>
</body>
</html>
\end{lstlisting}

Penjelasan tiap baris:

\begin{table}[ht]
\centering
\caption{Penjelasan struktur dasar HTML}
\label{tab:m03-baris}
\begin{tabular}{p{4cm}p{8cm}}
\toprule
\textbf{Baris} & \textbf{Fungsi} \\
\midrule
\texttt{<!DOCTYPE html>} & Memberitahu peramban bahwa ini dokumen HTML5 \\
\texttt{<html lang="id">} & Elemen root; \texttt{lang="id"} menandakan bahasa Indonesia \\
\texttt{<head>} & Metadata --- tidak terlihat di halaman, tapi penting untuk peramban \& SEO \\
\texttt{<meta charset>} & Set karakter \emph{encoding} (agar huruf seperti é, ü muncul benar) \\
\texttt{<meta name="viewport">} & Agar halaman responsif di \emph{mobile} \\
\texttt{<title>} & Judul yang muncul di tab peramban \\
\texttt{<body>} & Konten yang terlihat oleh pengguna \\
\bottomrule
\end{tabular}
\end{table}

% ---------------------------------------------------------------------
% 4. ELEMEN DASAR HTML
% ---------------------------------------------------------------------
\section{Elemen-Elemen Dasar HTML}
\label{sec:m03-elemen}

\begin{langkah}
  \item Tambahkan kode berikut di dalam \texttt{<body>}, di bawah \texttt{<h1>}.
\end{langkah}

\begin{lstlisting}[language=HTML, caption={Kumpulan elemen HTML dasar}, label={lst:m03-elemen}]
<!-- Heading: h1 sampai h6 -->
<h1>Ini Heading 1 (Judul Utama)</h1>
<h2>Ini Heading 2 (Judul Section)</h2>
<h3>Ini Heading 3 (Sub-judul)</h3>

<!-- Paragraph -->
<p>Ini adalah paragraf biasa. HTML menggunakan tag
   untuk mendefinisikan konten.</p>
<p>Paragraf kedua. Setiap tag &lt;p&gt; akan membuat
   baris baru.</p>

<!-- Link (Anchor) -->
<p>Kunjungi <a href="https://developer.mozilla.org"
   target="_blank">MDN Web Docs</a> untuk belajar
   HTML lebih lanjut.</p>

<!-- Image (gambar dari internet) -->
<img src="https://placehold.co/600x300"
     alt="Contoh gambar placeholder"
     width="600" height="300">

<!-- List tidak berurut (Unordered List) -->
<h3>Tools yang Kita Gunakan:</h3>
<ul>
    <li>VS Code --- editor kode</li>
    <li>Live Server --- preview otomatis</li>
    <li>Browser --- melihat hasil</li>
    <li>GitHub --- menyimpan kode</li>
</ul>

<!-- List berurut (Ordered List) -->
<h3>Langkah Instalasi:</h3>
<ol>
    <li>Download VS Code dari website resmi</li>
    <li>Install ekstensi Live Server</li>
    <li>Buat file index.html</li>
    <li>Klik kanan &rarr; Open with Live Server</li>
</ol>

<!-- Horizontal Rule -->
<hr>

<!-- Strong & Emphasis -->
<p>HTML itu <strong>sangat penting</strong> untuk
   dipelajari karena fondasi semua website.</p>
<p>Kamu pasti <em>bisa</em> menguasai ini!</p>
\end{lstlisting}

Simpan file (\textbf{Ctrl + S}) $\rightarrow$ buka dengan \emph{Live Server}.

\begin{tip}
Periksa di peramban: apakah terlihat 3 \emph{heading} (h1, h2, h3) dengan ukuran berbeda, 2 paragraf, 1 tautan, 1 gambar \emph{placeholder}, 1 daftar \emph{bullet}, 1 daftar bernomor, 1 garis horizontal, 1 teks tebal, dan 1 teks miring?
\end{tip}

\begin{tangkapanlayar}
  {assets/images/minggu-03/browser-elemen-dasar.png}
  {Peramban yang menampilkan hasil elemen HTML dasar.}
  {Pastikan terlihat heading, paragraph, link, gambar, lists, strong, em.}
\end{tangkapanlayar}

% ---------------------------------------------------------------------
% 5. SEMANTIC HTML
% ---------------------------------------------------------------------
\section{Semantic HTML}
\label{sec:m03-semantic}

\begin{catatan}
\textbf{Semantic HTML} adalah penggunaan tag yang tidak hanya mendeskripsikan tampilan, tapi juga \textbf{makna} konten. Ini penting untuk SEO, aksesibilitas, dan keterbacaan kode.
\end{catatan}

Perbedaan \emph{semantic} vs \emph{non-semantic}:

\begin{table}[ht]
\centering
\caption{Perbandingan tag semantic dan non-semantic}
\label{tab:m03-semantic}
\small
\begin{tabular}{p{3.5cm}p{3.5cm}p{5cm}}
\toprule
\textbf{Non-Semantic} & \textbf{Semantic} & \textbf{Kapan Menggunakan} \\
\midrule
\texttt{<div>} untuk semua \emph{section} & \texttt{<header>} untuk bagian atas & Gunakan tag semantik saat konten punya peran jelas \\
\texttt{<div>} untuk daftar menu & \texttt{<nav>} untuk navigasi & \texttt{<nav>} memberitahu peramban ini bagian navigasi \\
\texttt{<div>} untuk konten utama & \texttt{<main>} untuk konten utama & Hanya ada satu \texttt{<main>} per halaman \\
\texttt{<div>} untuk footer & \texttt{<footer>} untuk bagian bawah & Setiap halaman sebaiknya punya \texttt{<footer>} \\
\texttt{<div>} untuk \emph{section} & \texttt{<section>} atau \texttt{<article>} & \texttt{<section>} = bagian tematik; \texttt{<article>} = konten mandiri \\
\bottomrule
\end{tabular}
\end{table}

\subsection{Rekonstruksi dengan Semantic Tag}

Buat file baru bernama \texttt{profil.html} di folder yang sama. Ketik kode berikut:

\begin{lstlisting}[language=HTML, caption={Halaman profil dengan semantic HTML}, label={lst:m03-profil}]
<!DOCTYPE html>
<html lang="id">
<head>
    <meta charset="UTF-8">
    <meta name="viewport"
          content="width=device-width, initial-scale=1.0">
    <title>Profil Diri --- Latihan Semantic HTML</title>
</head>
<body>

    <!-- HEADER: Berisi navigasi -->
    <header>
        <h1>Portfolio Saya</h1>
        <nav>
            <ul>
                <li><a href="#tentang">Tentang</a></li>
                <li><a href="#keahlian">Keahlian</a></li>
                <li><a href="#proyek">Proyek</a></li>
                <li><a href="#kontak">Kontak</a></li>
            </ul>
        </nav>
    </header>

    <!-- MAIN: Konten utama halaman -->
    <main>

        <!-- SECTION: Tentang Saya -->
        <section id="tentang">
            <h2>Tentang Saya</h2>
            <img src="https://placehold.co/200x200"
                 alt="Foto profil"
                 width="200" height="200">
            <p>Halo! Nama saya [Nama Lengkap], mahasiswa
               semester 1 D3 Informatika di [Nama Kampus].
               Saya tertarik dengan dunia desain web dan
               ingin belajar membangun website yang menarik
               dan mudah digunakan.</p>
        </section>

        <!-- SECTION: Keahlian -->
        <section id="keahlian">
            <h2>Keahlian</h2>
            <ul>
                <li>HTML &amp; dasar CSS</li>
                <li>Figma (wireframing)</li>
                <li>Microsoft Office</li>
                <li>Dasar pemrograman (belajar)</li>
            </ul>
        </section>

        <!-- SECTION: Proyek -->
        <section id="proyek">
            <h2>Proyek</h2>

            <!-- ARTICLE: Konten mandiri tentang 1 proyek -->
            <article>
                <h3>Website Profil Diri</h3>
                <p>Website statis pertama saya yang dibuat
                   dengan HTML dan CSS. Berisi informasi
                   profil diri, daftar keahlian, dan cara
                   menghubungi saya.</p>
                <p><strong>Teknologi:</strong> HTML, CSS</p>
            </article>

            <article>
                <h3>Landing Page Produk</h3>
                <p>Halaman promosi produk sederhana yang
                   dibuat sebagai latihan layout dengan
                   Flexbox.</p>
                <p><strong>Teknologi:</strong> HTML, CSS</p>
            </article>
        </section>

        <!-- SECTION: Kontak -->
        <section id="kontak">
            <h2>Kontak</h2>
            <p>Email: <a href="mailto:nama@email.com">
               nama@email.com</a></p>
            <p>GitHub: <a href="https://github.com/username"
               target="_blank">github.com/username</a></p>
        </section>

    </main>

    <!-- FOOTER -->
    <footer>
        <p>&copy; 2026 [Nama Lengkap]. Dibuat untuk
           mata kuliah Desain Web.</p>
    </footer>

</body>
</html>
\end{lstlisting}

Simpan file $\rightarrow$ buka dengan \emph{Live Server}.

\subsection{Perbandingan dengan Non-Semantic}

\begin{lstlisting}[language=HTML, caption={Perbandingan non-semantic vs semantic}, label={lst:m03-perbandingan}]
<!-- Non-Semantic: semua pakai div -->
<div class="header">
    <div class="nav">
        <div class="menu-item">Tentang</div>
    </div>
</div>
<div class="main">
    <div class="section">...</div>
</div>
<div class="footer">...</div>

<!-- Semantic: menggunakan tag yang tepat -->
<header>
    <nav>
        <a href="#tentang">Tentang</a>
    </nav>
</header>
<main>
    <section>...</section>
</main>
<footer>...</footer>
\end{lstlisting}

\begin{catatan}
Kenapa \emph{semantic} lebih baik?
\begin{enumerate}
  \item \textbf{Screen reader} (untuk tunanetra) bisa membaca struktur halaman dengan benar.
  \item \textbf{Search engine} (Google) lebih mudah memahami konten halaman.
  \item \textbf{Developer lain} (dan kamu 3 bulan lagi) lebih mudah membaca kode.
\end{enumerate}
\end{catatan}

% ---------------------------------------------------------------------
% 6. ELEMEN PENDUKUNG PENTING
% ---------------------------------------------------------------------
\section{Elemen Pendukung Penting}
\label{sec:m03-pendukung}

\subsection{Anchor (Tautan) --- Detail Atribut}

\begin{lstlisting}[language=HTML, caption={Berbagai jenis tautan}, label={lst:m03-link}]
<!-- Link ke halaman lain -->
<a href="https://www.google.com">Google</a>

<!-- Link ke email -->
<a href="mailto:budi@email.com">Kirim Email</a>

<!-- Link ke section di halaman yang sama (anchor link) -->
<a href="#keahlian">Lihat Keahlian</a>

<!-- Buka di tab baru -->
<a href="https://www.google.com"
   target="_blank" rel="noopener noreferrer">
    Google (tab baru)
</a>

<!-- Gambar sebagai link -->
<a href="https://www.google.com">
    <img src="logo.png" alt="Logo Google" width="100">
</a>
\end{lstlisting}

\subsection{Image --- Detail Atribut}

\begin{lstlisting}[language=HTML, caption={Berbagai cara menampilkan gambar}, label={lst:m03-img}]
<!-- Gambar dasar -->
<img src="foto.jpg" alt="Deskripsi gambar"
     width="400" height="300">

<!-- Gambar dari URL -->
<img src="https://placehold.co/400x300" alt="Placeholder">

<!-- Gambar dengan figure & caption -->
<figure>
    <img src="foto.jpg" alt="Foto profil"
         width="200" height="200">
    <figcaption>Foto profil saya --- diambil tahun 2026
    </figcaption>
</figure>
\end{lstlisting}

\begin{peringatan}
Aturan penting untuk \texttt{alt} pada gambar:
\begin{itemize}
  \item \textbf{Selalu} isi atribut \texttt{alt}.
  \item \texttt{alt} harus mendeskripsikan gambar, bukan hanya ``gambar'' atau ``foto''.
  \item Jika gambar murni dekoratif, gunakan \texttt{alt=""} (kosong).
\end{itemize}
\end{peringatan}

\subsection{Tabel Dasar (Preview)}

\begin{lstlisting}[language=HTML, caption={Struktur tabel HTML sederhana}, label={lst:m03-tabel-preview}]
<table>
    <caption>Jadwal Kuliah Semester 1</caption>
    <thead>
        <tr>
            <th>Hari</th>
            <th>Mata Kuliah</th>
            <th>Jam</th>
        </tr>
    </thead>
    <tbody>
        <tr>
            <td>Senin</td>
            <td>Desain Web</td>
            <td>08:00 &ndash; 10:00</td>
        </tr>
        <tr>
            <td>Rabu</td>
            <td>Algoritma</td>
            <td>10:00 &ndash; 12:00</td>
        </tr>
    </tbody>
</table>
\end{lstlisting}

\begin{catatan}
Tabel akan dibahas lebih dalam pada Bab~\ref{chap:minggu-04}.
\end{catatan}

% ---------------------------------------------------------------------
% 7. MENGHUBUNGKAN HTML KE CSS
% ---------------------------------------------------------------------
\section{Menghubungkan HTML ke CSS (Dasar)}
\label{sec:m03-css}

\begin{catatan}
CSS detail akan dipelajari di minggu-minggu berikutnya. Ini hanya pengenalan cara menghubungkan file.
\end{catatan}

\subsection{Buat File CSS}

Di folder yang sama, buat file baru: \texttt{style.css}. Ketik kode berikut:

\begin{lstlisting}[language=CSS, caption={File CSS dasar untuk halaman profil}, label={lst:m03-css}]
/* Reset dasar */
body {
    font-family: Arial, sans-serif;
    margin: 0;
    padding: 0;
    line-height: 1.6;
    color: #333;
}

/* Header */
header {
    background-color: #2c3e50;
    color: white;
    padding: 20px;
    text-align: center;
}

header nav ul {
    list-style: none;
    padding: 0;
    display: flex;
    justify-content: center;
    gap: 20px;
}

header nav a {
    color: white;
    text-decoration: none;
}

/* Main content */
main {
    max-width: 800px;
    margin: 0 auto;
    padding: 20px;
}

section {
    margin-bottom: 40px;
}

/* Footer */
footer {
    background-color: #2c3e50;
    color: white;
    text-align: center;
    padding: 10px;
}
\end{lstlisting}

\subsection{Hubungkan ke HTML}

Di file \texttt{profil.html}, tambahkan tag \texttt{<link>} di dalam \texttt{<head>}:

\begin{lstlisting}[language=HTML, caption={Menambahkan tautan CSS di HTML}, label={lst:m03-link-css}]
<head>
    <meta charset="UTF-8">
    <meta name="viewport"
          content="width=device-width, initial-scale=1.0">
    <title>Profil Diri --- Latihan Semantic HTML</title>
    <link rel="stylesheet" href="style.css">
</head>
\end{lstlisting}

Simpan kedua file $\rightarrow$ buka \texttt{profil.html} via \emph{Live Server}.

\begin{tangkapanlayar}
  {assets/images/minggu-03/profil-dengan-css.png}
  {Halaman \texttt{profil.html} di peramban dengan CSS.}
  {Pastikan terlihat \emph{header} berwarna gelap, layout terpusat, dan \emph{footer} berwarna gelap.}
\end{tangkapanlayar}

% ---------------------------------------------------------------------
% 8. VALIDASI HTML
% ---------------------------------------------------------------------
\section{Validasi HTML}
\label{sec:m03-validasi}

\begin{langkah}
  \item Buka \texttt{https://validator.w3.org/\#validate\_by\_input}.
  \item Pilih tab \textbf{Direct Input}.
  \item Salin seluruh isi \texttt{profil.html} $\rightarrow$ tempel ke kolom teks.
  \item Klik \textbf{Check}.
  \item Pastikan hasil: ``Document checking completed. No errors or warnings to show.''
\end{langkah}

\begin{tangkapanlayar}
  {assets/images/minggu-03/w3c-validasi.png}
  {W3C HTML Validator yang menunjukkan hasil validasi tanpa error.}
  {Pastikan terlihat pesan ``No errors or warnings to show''.}
\end{tangkapanlayar}

Jika ada error, perbaiki berdasarkan pesan yang diberikan validator, lalu cek ulang.

% ---------------------------------------------------------------------
% 9. LATIHAN MANDIRI
% ---------------------------------------------------------------------
\section{Latihan Mandiri}
\label{sec:m03-latihan}

\begin{latihan}[Halaman ``Tips Belajar HTML'']
Buat file baru bernama \texttt{tips.html} dengan struktur \emph{semantic} HTML yang lengkap. Halaman ini berisi:
\begin{enumerate}
  \item \textbf{Header} dengan judul ``Tips Belajar HTML untuk Pemula'' dan \texttt{<nav>} (minimal 3 \emph{anchor link} ke \emph{section} di bawah).
  \item \textbf{Section ``Mengapa HTML Penting''} --- 2 paragraf penjelasan.
  \item \textbf{Section ``Tips Pertama''} --- \emph{ordered list} berisi minimal 5 tips.
  \item \textbf{Section ``Kesalahan Umum''} --- \emph{unordered list} berisi minimal 4 kesalahan, masing-masing dijelaskan dalam paragraf terpisah.
  \item \textbf{Section ``Sumber Belajar''} --- \emph{unordered list} berisi minimal 3 tautan ke website belajar HTML (MDN, W3Schools, dsb.). Setiap tautan harus menggunakan \texttt{target="\_blank"} dan \texttt{rel="noopener noreferrer"}.
  \item \textbf{Footer} dengan hak cipta tahun 2026.
\end{enumerate}
Persyaratan teknis: gunakan hanya tag \emph{semantic}; tidak boleh ada \texttt{<div>} sebagai \emph{section placeholder}; semua gambar harus punya \texttt{alt} yang deskriptif; file harus valid di W3C Validator.
\end{latihan}

\begin{latihan}[Push ke GitHub]
\begin{enumerate}
  \item Inisialisasi Git di folder \texttt{latihan-web-3} dan buat \emph{repository} baru di GitHub (public, dengan README).
  \item \emph{Push} kode:
\end{enumerate}
\end{latihan}

\begin{lstlisting}[caption={Perintah push ke GitHub}, label={lst:m03-git}]
git init
git remote add origin https://github.com/NAMA_USER/latihan-web-3.git
git add .
git commit -m "Latihan HTML minggu 3: semantic tags"
git branch -M main
git push -u origin main
\end{lstlisting}

\begin{tip}
Pastikan file \texttt{index.html}, \texttt{profil.html}, \texttt{tips.html}, dan \texttt{style.css} terlihat di \emph{repository} GitHub.
\end{tip}

% ---------------------------------------------------------------------
% 10. TROUBLESHOOTING
% ---------------------------------------------------------------------
\section{Troubleshooting}
\label{sec:m03-troubleshoot}

\begin{table}[ht]
\centering
\caption{Masalah umum HTML dan solusinya}
\label{tab:m03-troubleshoot}
\small
\begin{tabular}{p{4.5cm}p{8cm}}
\toprule
\textbf{Masalah} & \textbf{Solusi} \\
\midrule
Gambar tidak muncul di peramban & Periksa path \texttt{src} --- gunakan URL lengkap untuk gambar \emph{online}, atau pastikan nama file benar untuk gambar lokal \\
Tautan tidak berfungsi (\emph{anchor}) & Pastikan \texttt{href="\#section-id"} cocok dengan \texttt{id="section-id"} pada \emph{section} target \\
W3C Validator menampilkan error & Baca pesan error --- biasanya tag tidak ditutup (\texttt{</p>}) atau \emph{nesting} salah \\
CSS tidak terlihat & Pastikan \texttt{<link rel="stylesheet" href="style.css">} ada di \texttt{<head>}, dan nama file CSS benar \\
Halaman tidak responsif di \emph{mobile} & Tag \texttt{<meta name="viewport" ...>} harus ada di \texttt{<head>} \\
Teks berantakan / paragraf tidak terpisah & Pastikan setiap paragraf dibungkus tag \texttt{<p>}, bukan hanya enter di kode \\
\bottomrule
\end{tabular}
\end{table}

% ---------------------------------------------------------------------
% 11. RANGKUMAN
% ---------------------------------------------------------------------
\section{Rangkuman}
\label{sec:m03-rangkuman}

\begin{itemize}
  \item Struktur dasar dokumen HTML: \texttt{<!DOCTYPE html>}, \texttt{<html>}, \texttt{<head>}, \texttt{<body>}.
  \item Elemen dasar: \emph{heading}, paragraf, tautan, gambar, daftar, \texttt{<strong>}, \texttt{<em>}, \texttt{<hr>}.
  \item \emph{Semantic tags}: \texttt{<header>}, \texttt{<nav>}, \texttt{<main>}, \texttt{<section>}, \texttt{<article>}, \texttt{<footer>}.
  \item Perbedaan \emph{semantic} vs \emph{non-semantic} HTML.
  \item Menghubungkan HTML ke file CSS eksternal.
  \item Validasi HTML menggunakan W3C Validator.
\end{itemize}

\begin{tip}
\textbf{Cheat Sheet Tag Semantic:}
\begin{itemize}
  \item \texttt{<header>} $\rightarrow$ Bagian atas halaman (judul, \emph{nav})
  \item \texttt{<nav>} $\rightarrow$ Navigasi / menu
  \item \texttt{<main>} $\rightarrow$ Konten utama (hanya 1 per halaman)
  \item \texttt{<section>} $\rightarrow$ Bagian tematik dalam \texttt{<main>}
  \item \texttt{<article>} $\rightarrow$ Konten mandiri (bisa dipindah ke halaman lain)
  \item \texttt{<aside>} $\rightarrow$ Konten sampingan (\emph{sidebar})
  \item \texttt{<footer>} $\rightarrow$ Bagian bawah halaman
  \item \texttt{<figure>} $\rightarrow$ Gambar + \emph{caption}
  \item \texttt{<figcaption>} $\rightarrow$ \emph{Caption} untuk \texttt{<figure>}
\end{itemize}
\end{tip}
